% Implementation verification (Section S9). Moved from sections/verification.tex
% and sections/supplement.tex of the rule-assembled attempt; the main-text
% summary, whose every claim reappears in the record below, was merged into the
% overview during the journal revision (agent S, 2026-10-13). Test names and
% repository paths were replaced by a pointer to the archive README; claim
% markers (VERxx, AVxx, RCxx) refer to rows in claim-evidence.md.
\phantomsection\label{sec:verification-summary}\phantomsection\label{supp:verification-details}

Verification checks whether the implementation does what the frozen protocol specifies; it does not validate the social mechanism against data. Three layers checked the implementation before any confirmatory seed was used: unit tests, the procedural gates and a freeze preflight. Two docking studies then compared independent implementations with the project code. Three claims of different strength follow from this record, and they should not be confused. % VER01
\begin{itemize}
\item \emph{Intended determinism.} Execution is bound to the recorded environment: the freeze record binds the interpreter, NumPy, the lock file, the platform, single-threaded linear algebra and a numerical probe, and the runners refuse to start in any other environment. Within that environment and the frozen tree, the tests and outputs are designed to agree byte for byte; this is a property of the design, documented by no separate test-run record. % VER14 AV05
\item \emph{Demonstrated component agreement.} The docking studies showed agreement between independent implementations and the project code for specified components and cases, at reduced sizes, and a second host reproduced three G0 tasks to within $7.9\times10^{-12}$ in continuous values. These are checks of parts, not a reproduction of the release.
\item \emph{Unverified full-release replay.} Full-release bitwise replay has not been independently demonstrated, and the recorded checks on Linux failed their declared criteria. Since an operating-system update of 9 October 2026, after the run, the bound machine itself no longer passes the freeze check (Section~\ref{sec:verification-limits}). The release wrapper and the reporting loader also lie outside the frozen code closure (Section~\ref{sec:record-history}). % VER15
\end{itemize}
The README of the archived code lists every test by the property it checks and gives the commands of the runners, the preflight and the reporting steps. Table~\ref{tab:audit-map} (Section~\ref{sec:rule-record}) maps each kind of claim to its public evidence, to what was not kept and to what was checked independently.

\subsection{Test suite}

% VER02
The unit tests belong to the frozen code closure and are grouped here by the property they check; the archive README names them.
\begin{itemize}
\item \emph{Identities (the identity unit tests, G1a).} At full $N$, the zero-strength influence arm and persistence models A, B, D and E must coincide byte for byte. This covers the mask, measured, observed and selected features, in the matching and validation phases, on both attempts, and across the main, mobility, persistent-style, conditioning and 24-step variants. The optimized transport must reproduce the earlier generator's bytes, and the selected features must be byte-identical between implementations. An anchor-exact fit must give zero shifts and zero nuisance variance.
\item \emph{Persistence laws.} Tests check the declared transported laws of B and E, and the degeneracy at $\omega=1$, which is independent of strength and carries no identity claim.
\item \emph{Profile solver.} Tests cover roots on every face and corner; impossible targets, non-finite values and exact budget counting; non-monotone and flat profiles; degenerate secants; two historical stall regressions; the quarter-standard-error acceptance rule; the budget of 96 actual evaluations; and the rule that an endpoint reached without a solver step is labeled \code{boundary}.
\item \emph{Validation and retry.} Tests check the combined-variance validation rule and check that a validation failure leads to a single retry, which keeps both attempts. That the retry happens only after a validation failure and has its own budget is tested only through two pre-freeze scripts, one that checks recorded matching measurements and one that projects the compute cost. Neither property is tested on the fitting code, in which each attempt receives its own budget of 96 evaluations by construction.
\item \emph{Nuisance propagation.} The mean-feature shift predicted at a fixed test point must match regenerated panels within 5\% at the vector level. Further tests cover projection with lag-1 priority, a centered variance with cluster resampling, and uncentered boundary draws.
\item \emph{Band, censoring and scheduling.} Tests cover the nine-strength band, the zero-strength check when the cell at $a=0$ is available (no test checks the rule with that cell unavailable), censoring only across gaps, the survival of other units' boundaries when one unit is invalid, the fixed tier order with its barriers, and the frozen runner's refusal to release before the original deadline. These tests describe the frozen runner, while the linked release amendment is implemented by a later wrapper.
\item \emph{Gate criteria and seeds.} Tests check that failures stay in the denominators, that the finite-difference check uses one median per cell, that the pooled allowance scales with the allocated cells, that the nuisance-variance check reuses the baseline seeds, and that seed domains do not alias, with the confirmatory namespaces reserved. A separate seed-role audit record enumerates seed materials across gate plans and namespaces to check for collisions.
\end{itemize}
% VER14
The test suite depends on the environment, so byte-level checks are expected to hold only in the bound environment; the Linux failures of the frozen-calibration byte test are described under the environment diagnostics below.

\subsection{Environment diagnostics and reproduction limits}\label{sec:verification-limits}

Bitwise agreement is intended only within the bound environment and the frozen tree; no separate test-run record documents it, and complete-run replay has not been independently demonstrated. Later revisions of the repository change the dependency declaration and the lock file, which are in the bound code closure, and add Python files that the preflight also hashes, so the preflight fails on them. % AV05
On Linux x86-64 (Python 3.12.3, NumPy 2.5.3 from the lock), 538 study tests pass and 3 fail, and the same three failures reproduce in a clean worktree of revision \texttt{37513841}: the two cases of the frozen-calibration byte-hash test ($N=600$ and $N=10{,}626$) and one profile-solver test. % AV06
The cause has not been verified; the working hypothesis is that the recorded hashes were generated on macOS and that floating-point arithmetic differs between platforms. % AV06
A second host (Linux x86-64, with a different build of Python 3.12.13) replayed three G0 tasks for equivalence only, not as gate evidence. % AV07
All 82 code hashes matched the freeze record, the numerical probe agreed within $6.7\times10^{-16}$, states, decisions and counts were equal, and 1,331,544 continuous values agreed to within $7.9\times10^{-12}$ in absolute terms. % AV07
The pre-declared equivalence criterion nevertheless failed, only on the \code{seed_material_sha256} field of the phase records (4,448 divergences), because that hash is computed over the numerical parameters at which panels are generated; the platform, interpreter and host fields of the run manifests also differ, as expected between hosts. % AV08
Gate stages and tiers were therefore run on the one bound machine. % AV09
On 9 October 2026, after the run, an operating-system update of that machine (macOS 27.0 to 27.0.1) changed its platform string from the value bound in the freeze record, so the frozen preflight would now refuse to run the gate or the tiers there; the reporting steps do not call the preflight (decision log item 270). The bound environment is thus recorded, but the bound machine itself no longer passes the freeze check.

\paragraph{Replication route.}
Every result table and figure of the article except Table~\ref{main-tab:exploratory}, and every table and figure of this resource rendered by the journal export, regenerate from the archived aggregates with the archived scripts. Table~\ref{main-tab:exploratory} is written by hand, with each result number taken from the archived values files. The tables of this resource rendered by the report export are edited copies of report tables that the archived scripts regenerate, and their result numbers keep their value keys. The other tables of Section~\ref{sec:exploratory} were rendered by the exploratory export, which needs the aggregates of the exploratory analyses X3, X4, X6 and X7; these are not archived because they contain host and path records, so those tables cannot be regenerated from the archive. Their values are in the archived values catalog, which is why, unlike the twin, variance and stress-test aggregates, no sanitized copies of them were made. The values of Tables~\ref{tab:supp-x1a} and~\ref{tab:supp-x2} were transcribed without a script from the archived exploratory aggregates (Section~\ref{sec:rule-record}). Re-running the simulation outside the bound environment, including on the bound machine since its update, means calling the library entry points of the frozen code without the runners' freeze check; the archive README describes this route. The outputs are then not bitwise comparable with the released ones. The random seeds are derived from text labels and do not change, but floating-point results can differ in the last digits (in the cross-host replay above, continuous values agreed to within $7.9\times10^{-12}$); record keys that hash the numerical parameters, such as \code{seed_material_sha256}, then differ, and a decision that falls at a tolerance edge can differ too. Agreement of a complete re-run outside the bound environment with the released results, statistically or within the reported bands, has not been demonstrated and must be assessed rather than assumed.

\paragraph{Correction to the exploratory-aggregate README.}
The earlier statement in \code{archive/exploratory/README.md} that every table and figure regenerates except the twin and variance tables is too broad for this journal version. The applicable scope is given above: article Table~\ref{main-tab:exploratory} is written by hand, and the supplementary tables rendered by the exploratory export cannot be regenerated from the archive because the necessary aggregates are not archived. Their reported values remain available as described above. This correction supersedes that statement without modifying the preserved sanitized aggregates or their original README.

\subsection{Gates and preflight}

% VER03
The validation gate (Section~\ref{sec:odd}, Submodel~14) checks the procedure end to end; its stages, criteria and outcomes are in Section~\ref{sec:record-gate}. The release wrapper and the reporting loader lie outside the frozen code closure, and the repaired summary of the zero-strength band check is stored outside the seal (Section~\ref{sec:record-history}); the actual release time is held in a separate release record (Section~\ref{sec:record-execution}).

% VER04
Before any gate stage or tier runs, the preflight compares the working tree with the freeze record. It requires:
\begin{itemize}
\item schema 2 and status \code{frozen};
\item SHA-256 equality for every file in the bound code closure, including the lock file, and for the criteria and inventory artifacts;
\item agreement of the criteria record with the criteria code and calibration script;
\item hashes of the amendment, specification, revision, review findings and both decision records;
\item eight workers, eight shards and single-threaded BLAS;
\item the recorded interpreter, NumPy, platform and lock environment;
\item agreement of a numerical probe within absolute tolerance $10^{-12}$.
\end{itemize}
A second binding record, created after G0 and before G2, must also exist. It binds the freeze record, the G0 statuses, the criteria record, the G2 cell allocation and the measured G0 gap rate, and is checked against them before G2. Tests check that the runners refuse to start without these records or with stale or changed inputs.

\subsection{Independent docking}\label{sec:verification-docking}

% VER06
Two clean-room docking studies compared independent implementations with the project code at the frozen revision \texttt{37513841}. Each independent implementation was written from the protocol and specification documents alone, without opening the project's source, tests or experiment scripts. Its first-phase readings of ambiguous points were committed before any project source was read, and readings added after the comparison are marked as such. A comparison plan was committed before any comparison ran. In the model study, a timing probe had earlier printed mean lag correlations of 50 project panels in two cells, which Section~\ref{sec:record-exposure} lists as exposure; no margin or cell was chosen from it. Neither study ran a gate stage or tier or used a reserved confirmatory seed namespace. The phase commits in the repository were reconstructed from the session transcripts in their original order; the original commit objects were not pushed, which limits the strength of a clean-room claim. The docking reports and reconstruction notes are in the archive.

% VER07
The first study docked the generative model. With shared seeds, latent trajectories were bitwise identical for the influence arm at $a\in\{0,0.03,0.30\}$ and for persistence models A, B, D and E, in 20 of 20 seeds per case. In a post hoc check not in the pre-declared plan, the ordinal discretizers were also byte-identical given the same continuous responses and mask. At the level of laws ($N=1{,}002$, 11 waves, 25 cells, 1,000 panels per implementation and cell), all 150 primary endpoints were equivalent by two one-sided tests with a margin of 0.20 standard deviations per panel. None of 2,334 secondary difference tests rejected after Holm correction. For the discriminator, run at reduced training and test sizes in three cells with 50 replicates each, mean AUCs differed by at most 0.002 (equivalence margin 0.02). On identical features, the two LDA implementations differed by at most $9\times10^{-5}$. These reduced-size AUCs are not estimates of $\mathcal A(a,r)$.

% VER08
Two specification ambiguities affected numerical output, both negligibly. The first concerned the LDA standardization divisor (the project uses the population standard deviation). The second concerned the normalization of the heteroscedastic error in C (the project uses the empirical mean over the panel, which differed from the analytic expectation by a relative $1.6\times10^{-4}$). No discrepancy in the docked scope was classified as a code--specification divergence. The study did not dock matching, nuisance propagation, bootstrap bands, gates, seeds or sealing; the latent, error-free ordinal, latent-rank and mobility channels; the persistent-style, conditioning, anchoring, attraction-only and 24-step sensitivities; the lag-one feasibility rule; or production sizes. Its law-level resolution (0.20 standard deviations, about 0.002--0.005 in units of the ordinal lag correlations) is far coarser than the matching tolerance. Heaping and masks agree only in law because the two implementations consume random draws in a different order.

% VER09
The second study docked the inference and calibration machinery: the profile solver, acceptance tolerance, independent validation and retry, local nuisance propagation, cluster bootstrap, band, gap and gate rules, and the Jeffreys stop rule. Both implementations received identical synthetic inputs, including adversarial ones. Tolerances, validation decisions, the retry protocol, projected parameter shifts (2,500 cases), finite differences, nuisance draws, $V_{\mathrm{nuis}}$, bootstrap standard errors, failure allowances and the shared-seed stop-rule probability agreed to $10^{-12}$ or better, most bitwise. Exceptions were reporting-only differences at 0/0 and non-finite inputs and in one Monte Carlo standard-error formula. At a non-positive Schur slope, the project stops with an error where the docking solves the projection; this was classified as a specification ambiguity. For the nuisance draws, the docking implementation was given the project's discriminator.

% VER10
Solver status agreed in 16 of 22 synthetic cases; the report's count of 17 is corrected in the study's reconstruction note. All six divergences (cases 5, 6, 9, 10, 17 and 21 of the docking record) were traced to four specification ambiguities, which that record lists. In three cases the project reports \code{no-L1-root} where the docking reports a family limit; both labels are gaps. The project can also report family limits within the profile solver, so these cases do not imply that such limits arise only in the feasibility screen. In one case with discrete steps the two labels are \code{numerical-limit} and \code{ill-conditioned}. In one case a root lies within tolerance of the endpoint and is labeled \code{interior} by the project and \code{boundary} by the docking; both are accepted states. In the hump-shaped profile case, the project accepted a genuine interior root after sampling only $\eta\leq0.01$, whereas the docking evaluated $\eta=0.8$ early, detected the non-monotone profile and returned a gap. This is the only divergence that changes the accept-or-gap decision.

% VER11
The band critical values differed, by up to a factor of 2.1, in the ten random cases compared. A post hoc analysis reproduced the project's values exactly in all ten cases from three readings: the logit standard error evaluated at each replicate rather than at $\hat p$, replicates centered on the point estimate, and the \texttt{higher} quantile rule. The per-replicate standard error dominates; in one case the plug-in reading gave a critical value of 7.15 and the project's reading 3.37. The project's band code reproduces the frozen v1 band exactly, and Section~\ref{sec:design-details} describes that reading. Inversion, gap, censoring and zero-strength rules agreed in 9 of 11 cases. In the other two, adjacent to a gap, the project censors the sustained boundary where the docking declares the first strength above the gap, which is the more conservative reading. No divergence was classified as a probable project error.

% VER12
The second study left matching on real panels, the feasibility screen and the mean-feature Jacobian on real realizations unverified. It also left unverified the control path that inherits $\eta$ in the nuisance draws; the amplification and covariance diagnostics; and the calibrated G2 dispersion thresholds and the G2c ratio beyond failure counts. G1b, G1c, the G0 status projection that feeds the stop rule, the curve runner, scheduler, sealing, seed derivation and collision audit also remained unverified. The study did not verify solver efficiency on real targets, since evaluation counts were reported but not compared, or the statistical coverage of either band reading. The project's reading is the one that the band calibration (G1b) tests.

% VER13
Docking establishes agreement between two implementations of the same written specification in the scope docked. Where both made the same choice at an ambiguous point, it cannot detect an error in the specification itself or a misreading shared by both authors. The two implementations are also less independent than those of two human teams: the frozen code was implemented with AI coding agents (Codex and Claude) and the docking implementations were written by Claude (Section~\ref{sec:ai}), so a misreading shared by those models would go undetected. Docking is not validation of the social mechanism, not an estimate of coverage and not a substitute for the gates.
